\chapter*{Contribuciones Personales}
\label{cap:contribucionesPersonales}
\addcontentsline{toc}{chapter}{Contribuciones Personales}

A lo largo de este Trabajo de Fin de Grado hemos trabajado de forma coordinada, distribuyendo las responsabilidades de manera que cada integrante del equipo asumiera un ámbito técnico principal sin perder la visión global del proyecto. En las secciones siguientes se detallan las contribuciones individuales de cada uno.

\section*{Contribuciones de Sheila Julvez López}

\subsection*{Investigación del estado de la cuestión y selección tecnológica}

La investigación del estado del arte fue liderada principalmente por Sheila. Esta fase, que se corresponde con el capítulo~\ref{cap:estadoDeLaCuestion} de la memoria, implicó revisar en profundidad la literatura académica y la documentación técnica sobre realidad aumentada: sus orígenes militares en los HUDs, su evolución hasta los dispositivos de consumo actuales, la taxonomía de tipos de RA (basada en marcadores, sin marcadores, óptica, proyectiva), y las aplicaciones en entornos de retransmisión audiovisual y \textit{broadcasting} profesional.

Sheila analizó y comparó los principales SDKs de realidad aumentada disponibles en el mercado: Vuforia, ARToolKit, los SDK de NVIDIA, DeepAR, ARKit y ARCore. Para cada uno evaluó la compatibilidad con modelos 3D, la facilidad de integración como plugin nativo, el rendimiento en tiempo real y la política de licencias. Esta revisión comparativa fue determinante para justificar la elección de OpenCV y la biblioteca Assimp como base tecnológica del proyecto, descartando soluciones que habrían requerido procesos externos o modificaciones del núcleo de OBS Studio.

Además, investigó a fondo la arquitectura interna de OBS Studio: su sistema de escenas y fuentes, el modelo de filtros de vídeo y el ciclo de vida de un \textit{plugin} nativo. Estudió la API pública de OBS y los \textit{templates} de desarrollo disponibles en el repositorio oficial, sentando las bases para que la integración del \textit{plugin} pudiera hacerse respetando la arquitectura nativa del software.

\subsection*{Desarrollo de la herramienta de calibración de cámara}

Una de las contribuciones técnicas individuales más relevantes de Sheila fue el diseño e implementación de la herramienta de calibración de cámara, desarrollada como utilidad independiente en Python. Esta herramienta automatiza el proceso de adquisición de imágenes de un tablero de ajedrez desde la cámara conectada, extrae las esquinas de los cuadros con subpíxel de precisión, y resuelve el modelo de cámara pinhole para obtener los parámetros intrínsecos (distancia focal en píxeles, punto principal) y los coeficientes de distorsión radial y tangencial.

El fichero de calibración resultante en formato XML es cargado por el \textit{plugin} al arrancar y aplicado a cada fotograma antes de la estimación de pose, corrigiendo las deformaciones ópticas propias de cada objetivo. Sheila realizó varias sesiones de calibración con distintas cámaras para validar la estabilidad y repetibilidad de los parámetros obtenidos, y documentó el procedimiento de uso de la herramienta tanto en el código fuente como en el apéndice correspondiente de esta memoria.

\subsection*{Diseño de la interfaz de usuario en OBS}

Sheila fue la principal responsable del diseño y la implementación de la interfaz de usuario del \textit{plugin} dentro de OBS Studio. Definió los controles de propiedad que el operador puede ajustar en directo desde el panel de filtros: rutas a los ficheros de calibración y de modelo 3D, la URL de la API de DOMjudge, el modo de operación activo y los parámetros de escala y posición del contenido superpuesto.

El diseño de la interfaz buscó que cualquier operador pudiera configurar y cambiar los parámetros del \textit{plugin} sin necesidad de reiniciar la escena ni conocer los detalles internos del sistema. Sheila realizó varias pruebas de usabilidad con usuarios con distinto nivel técnico, incorporando mejoras iterativas para que los controles resultaran intuitivos incluso en condiciones de producción con poco tiempo disponible.

\subsection*{Pruebas, validación y documentación}

Sheila coordinó la fase de pruebas y validación del sistema. Diseñó una batería de pruebas funcionales que cubría los cuatro modos de operación del \textit{plugin} bajo distintas condiciones: iluminación variable, distintas distancias al marcador, diferentes resoluciones de cámara y distintas tasas de fotogramas configuradas en OBS. Registró sistemáticamente los resultados y colaboró con Jose en la corrección de los problemas detectados.

En cuanto a la documentación, Sheila redactó los capítulos introductorios y el capítulo del estado de la cuestión, así como las secciones de motivación, objetivos y conclusiones de la memoria. También se encargó de la revisión de estilo y coherencia del conjunto del documento, unificando la terminología y la estructura de las referencias bibliográficas.

\section*{Contribuciones de Jose Moreno Barbero}

\subsection*{Arquitectura del \textit{plugin} y ciclo de renderizado}

Jose fue el principal responsable del diseño e implementación de la arquitectura central del \textit{plugin}. Esto implicó, en primer lugar, estudiar en detalle el modelo de extensión de OBS Studio y comprender el ciclo de vida de un filtro de vídeo nativo: los puntos de entrada \texttt{video\_render} y \texttt{video\_tick}, la gestión de la memoria de \textit{frame}, y las restricciones de tiempo que impone el ciclo de codificación de OBS para no introducir latencia adicional en la señal de salida.

A partir de ese análisis, Jose diseñó la estructura modular del \textit{plugin}, separando claramente el módulo de visión artificial (detección de marcadores y estimación de pose), el módulo de renderizado 3D y el módulo de conectividad con la API externa. Esta separación de responsabilidades permite añadir nuevos modos de operación sin modificar el código de los módulos existentes y facilita el mantenimiento a largo plazo.

\subsection*{Detección de marcadores ArUco y estimación de pose}

Jose implementó el subsistema de detección de marcadores ArUco mediante la biblioteca OpenCV. El módulo convierte cada fotograma recibido desde OBS al espacio de color adecuado para el detector (lo que requirió implementar un conversor que soporta los siete formatos de píxel que OBS puede entregar: I420, NV12, YUY2, UYVY, BGRA, RGBA y Y800) y ejecuta la detección y la estimación de pose con el algoritmo PnP sobre los parámetros de calibración cargados al arrancar.

Uno de los problemas técnicos más complejos que Jose resolvió fue la discrepancia entre los sistemas de coordenadas de OpenCV y OBS. OpenCV usa una convención con el eje Y apuntando hacia abajo y el eje Z hacia el frente de la cámara, mientras que OBS opera con una convención distinta. La solución requirió aplicar una rotación de $180^\circ$ sobre el eje horizontal y negar las componentes $Y$ y $Z$ del vector de rotación ArUco, transformación que Jose derivó analíticamente y validó empíricamente comparando la orientación renderizada con la real.

Además, Jose implementó el filtro de media exponencial con interpolación circular para suavizar la estimación de pose fotograma a fotograma. El tratamiento especial del salto de ángulo en la interpolación circular (necesario para evitar que el suavizado produzca rotaciones incoherentes al cruzar la discontinuidad de $\pm\pi$) fue una de las aportaciones más delicadas del sistema, ya que un error en este punto genera un temblor visual imposible de corregir en postproducción.

\subsection*{Renderizado 3D con Assimp y modos de operación}

Jose implementó el subsistema de renderizado 3D, que utiliza la biblioteca Assimp para cargar modelos en múltiples formatos (OBJ, FBX, GLTF, entre otros) y los dibuja sobre el fotograma de OBS con la posición y orientación estimadas por el módulo de pose. El renderizado calcula la proyección perspectiva a partir de los parámetros intrínsecos de la cámara, lo que garantiza que el modelo virtual aparezca correctamente escalado y situado independientemente de la distancia al marcador.

A partir de esta base, Jose desarrolló los cuatro modos de operación del \textit{plugin}. En el \textbf{modo de modelo 3D}, el sistema superpone geometría arbitraria anclada al marcador. En el \textbf{modo cronómetro}, el modelo actúa como reloj analógico de cuenta atrás cuyas manecillas se actualizan fotograma a fotograma y se sincronizan periódicamente con el tiempo oficial del concurso obtenido desde la API de DOMjudge. En el \textbf{modo \textit{scoreboard}}, el sistema proyecta la clasificación completa del concurso sobre el vídeo en tiempo real. En el \textbf{modo de información de equipo}, el sistema identifica automáticamente qué equipo tiene la cámara en el campo de visión y muestra exclusivamente sus datos, combinando la detección de marcadores con un fichero JSON de mapeo preparado antes de la retransmisión.

\subsection*{Integración con DOMjudge y arquitectura de dos hilos}

Jose diseñó e implementó la capa de conectividad con la plataforma de juez automático DOMjudge. El cliente de red consume la API REST de DOMjudge usando libcurl, biblioteca HTTP incluida en OBS Studio, para obtener de forma periódica la clasificación general del concurso, los veredictos de los envíos y el tiempo oficial restante.

Para que las peticiones HTTP no bloqueen en ningún caso el ciclo de renderizado, Jose diseñó una arquitectura de dos hilos: el hilo principal ejecuta la detección de marcadores, la estimación de pose y el renderizado dentro del \textit{callback} de OBS, mientras que un hilo secundario realiza las peticiones de red de forma asíncrona. El acceso concurrente a los datos compartidos se protege mediante una sección crítica (\texttt{CRITICAL\_SECTION}), garantizando que la actualización de datos nunca interrumpa la producción de fotogramas. Esta arquitectura fue diseñada explícitamente para entornos de producción en directo, donde cualquier parada del renderizado se traduce en un fotograma perdido perceptible por el espectador.

\subsection*{Capítulos técnicos de la memoria}

Jose redactó los capítulos de descripción técnica del trabajo: el capítulo~\ref{cap:descripcionTrabajo}, que detalla el diseño e implementación del filtro de RA base, y el capítulo~\ref{cap:descripcionTrabajoConcursosProgramacion}, que describe la extensión del sistema para los concursos de programación. Ambos capítulos incluyen diagramas de arquitectura, fragmentos de código comentados y análisis de las decisiones de diseño tomadas en cada etapa. Jose también redactó los apéndices técnicos y la sección de trabajo futuro de las conclusiones.
